\clearpage
\section*{Библиография}
\addcontentsline{toc}{section}{Библиография}
\label{sec:bibliografiya}

% Списъкът е с разстояние 1,5 реда като основния текст; размерът на шрифта е 11 pt.
{\small

\subsection*{Литература}
\begin{biblist}[series=bib]
\item Владова-Недкова, И. Разследване на компютърни престъпления. София: Софи-Р, 2023.
\item Вучков, В. Електронните доказателства по наказателни дела и техните източници~--
  необходими разсъждения.~-- Юридическо списание на НБУ, 2023, т.~19, №~2, с.~12--19.
\item Вучков, В. Наказателен процес на Република България. София: Сиби, 2026.
\item Вучков, В. Образуване на досъдебно производство: актуални проблеми. София: Сиби,
  2015.
\item Иванова, И. Конституционни проблеми във връзка с правомощията на прокурора по
  чл.~145 от Закона за съдебната власт.~-- Юридически свят, 2023, №~2, с.~32--59.
\item Манев, А. Ролята на прокурора при привличането на обвиняемия и повдигането на
  предварителното обвинение.~-- Общество и право, 2023, №~8, с.~44--60.
\item Манев, Н. Наказателнопроцесуално право: ръководство за студенти. София: Сиела, 2016.
\item Младенов, Д. Електронни доказателства.~-- De Jure, 2022, т.~13, №~1, с.~63--74.
\item Николова, Л. Прокурорът в досъдебното производство според новия
  Наказателно-процесуален кодекс.~-- Адвокатски преглед, 2006, №~2--3, с.~4--22.
\item Чинова, М. Досъдебното производство по НПК: теория и практика. София: Сиела, 2013.
\item Чинова, М., Г.~Митов. Кратък лекционен курс по наказателнопроцесуално право. София:
  Сиела, 2021.
\item Council of Europe. Electronic Evidence Guide. Version~2.1. Strasbourg: Council of
  Europe, 2020.
\item Council of Europe. Explanatory Report to the Convention on Cybercrime. European
  Treaty Series No.~185. Budapest, 23.11.2001.
\end{biblist}

\subsection*{Нормативни актове}
\begin{biblist}[resume*=bib]
\item Конституция на Република България. Обн. ДВ, бр.~56 от 13.07.1991~г.; последно изм.
  ДВ, бр.~66 от 06.08.2024~г.
\item Наказателно-процесуален кодекс. Обн. ДВ, бр.~86 от 28.10.2005~г., в сила от
  29.04.2006~г.; последно изм. ДВ, бр.~59 от 30.06.2026~г.
\item Наказателен кодекс. Обн. ДВ, бр.~26 от 02.04.1968~г.; последно изм. ДВ, бр.~55 от
  16.06.2026~г.
\item Закон за съдебната власт. Обн. ДВ, бр.~64 от 07.08.2007~г.; последно изм. ДВ,
  бр.~99 от 30.10.2026~г.
\item Закон за електронните съобщения. Обн. ДВ, бр.~41 от 22.05.2007~г.; последно изм.
  ДВ, бр.~66 от 21.07.2026~г.
\item Закон за европейската заповед за разследване. Обн. ДВ, бр.~16 от 20.02.2018~г.;
  последно изм. ДВ, бр.~36 от 2024~г.
\end{biblist}

\subsection*{Международни договори и актове на Европейския съюз}
\begin{biblist}[resume*=bib]
\item Конвенция за престъпления в кибернетичното пространство. Будапеща, 23.11.2001~г.
  Ратифицирана със закон, обн. ДВ, бр.~29 от 05.04.2005~г.; обн. ДВ, бр.~76 от
  15.09.2006~г.; в сила за Република България от 01.08.2005~г.
\item Second Additional Protocol to the Convention on Cybercrime on enhanced co-operation and
  disclosure of electronic evidence. Council of Europe Treaty Series No.~224. Strasbourg,
  12.05.2022.
\item Регламент (ЕС) 2017/1939 на Съвета от 12 октомври 2017~г. за установяване на засилено
  сътрудничество за създаване на Европейска прокуратура. ОВ L~283, 31.10.2017~г., с.~1.
\item Регламент (ЕС) 2023/1543 на Европейския парламент и на Съвета от 12 юли 2023 година
  относно европейските заповеди за предоставяне и европейските заповеди за запазване на
  електронни доказателства в рамките на наказателните производства и за изпълнението на
  наказания лишаване от свобода вследствие на наказателни производства. ОВ L~191,
  28.7.2023~г., с.~118--180.
\item Директива (ЕС) 2023/1544 на Европейския парламент и на Съвета от 12 юли 2023 година
  за установяване на хармонизирани правила относно определянето на определени форми на
  установяване и определянето на представители за целите на събирането на електронни
  доказателства по наказателни производства. ОВ L~191, 28.7.2023~г., с.~181.
\end{biblist}

\subsection*{Съдебна практика}
\begin{biblist}[resume*=bib]
\item Решение №~2 от 12.03.2015~г. на Конституционния съд по к.д. №~8/2014~г. (обн. ДВ,
  бр.~23 от 2015~г.).
\item Решение №~14 от 26.07.2024~г. на Конституционния съд по к.д. №~10/2023~г.
\item Решение на Съда на ЕС от 2 март 2021~г. по дело C-746/18, \textit{Prokuratuur},
  ECLI:EU:C:2021:152.
\item Решение на Съда на ЕС от 16 декември 2021~г. по дело C-724/19,
  ECLI:EU:C:2021:1020.
\item Решение на Съда на ЕС от 17 ноември 2022~г. по дело C-350/21, \textit{Spetsializirana
  prokuratura}, ECLI:EU:C:2022:896.
\item Решение на Съда на ЕС от 4 октомври 2024~г. по дело C-548/21,
  \textit{Bezirkshauptmannschaft Landeck}, ECLI:EU:C:2024:830.
\item Решение на ЕСПЧ от 5 ноември 2009~г. по делото „Колеви срещу България“, жалба №~1108/02.
\item Решение на ЕСПЧ от 11 януари 2022~г. по делото „Екимджиев и други срещу България“, жалба №~70078/12.
\end{biblist}

\subsection*{Други източници}
\begin{biblist}[resume*=bib]
\item Проект на Закон за изменение и допълнение на Наказателно-процесуалния кодекс
  (Министерство на правосъдието). Обществена консултация №~12669-K,
  \url{https://www.strategy.bg/bg/public-consultations/12669}.
\end{biblist}

}
